\title{DeepSeekMath: Pushing the Limits of Mathematical Reasoning in Open Language Models}

\begin{abstract}
Language models frequently struggle with mathematical reasoning, owing to the intricate and highly structured demands of mathematical tasks. This study presents DeepSeekMath~7B, developed by extending the pre-training of DeepSeek-Coder-Base-v1.5 7B using 120B tokens with a strong emphasis on mathematics, extracted from Common Crawl and supplemented with additional natural language and code resources. Without employing external tools or ensemble strategies, DeepSeekMath~7B attains a remarkable 51.7\% on the MATH benchmark at the competition level, nearing the accuracy levels of models such as Gemini-Ultra and GPT-4. Furthermore, when aggregating 64 samples, self-consistency elevates its performance to 60.9\% on MATH. The notable improvement in mathematical reasoning for DeepSeekMath~7B arises from two central innovations: first, a carefully crafted data selection pipeline maximizes the utility of large-scale public web datasets; second, the adoption of Group Relative Policy Optimization (GRPO), an adaptation of Proximal Policy Optimization (PPO), advances both mathematical reasoning aptitude and the computational efficiency of PPO.
\end{abstract}

\section{Introduction}

Large language models (LLMs) have transformed mathematical reasoning within artificial intelligence, leading to major progress on benchmarks for quantitative reasoning \citep{MATH} as well as for geometry reasoning \citep{trinh2024solving}. In addition, LLMs have demonstrated their value as tools that aid humans in tackling challenging mathematical tasks \citep{tao}. Despite these successes, state-of-the-art systems such as GPT-4 \citep{gpt4} and Gemini-Ultra \citep{gemini} remain inaccessible to the general public, while current open-source alternatives lag significantly in terms of performance.

This paper presents DeepSeekMath, a language model tailored to mathematical domains that demonstrates considerable advancements over existing open-source models, coming close to matching GPT-4’s performance on standard academic benchmarks. Central to our approach is the construction of the DeepSeekMath Corpus—a large-scale, high-quality pre-training dataset containing 120B math tokens. The data for this corpus is sourced from the Common Crawl (CC) using a classifier built on fastText \citep{joulin2016fasttext}. In its initial phase, the classifier is trained with positive examples from OpenWebMath \citep{openwebmath} and a varied set of unrelated web pages as negative samples. The trained classifier is then used to identify further math-related documents from the CC. These results are curated with the aid of human annotators and incorporated into the training set to continually refine the classifier's accuracy. Evaluation demonstrates the efficacy of our pre-training corpus: the base variant, DeepSeekMath-Base 7B, achieves 64.2\% on GSM8K \citep{gsm8k} and 36.2\% on the challenging MATH benchmark \citep{MATH}, surpassing results from Minerva 540B \citep{minerva}. Additionally, the multilingual nature of the DeepSeekMath Corpus results in noticeable gains on Chinese mathematical benchmarks \citep{wei2023cmath,agieval}. We contend that our methods for collecting and processing mathematical data offer a valuable foundation for further research, with substantial opportunities for continued progress.

DeepSeekMath-Base is built upon DeepSeek-Coder-Base-v1.5 7B \citep{deepseek-coder}, since we find that models pretrained on code exhibit stronger performance than those beginning from standard large language models. Additionally, our results show that training on mathematical tasks also yields improvements in model performance on the MMLU \citep{mmlu} and BBH \citep{bbh} benchmarks. This suggests that mathematical training not only strengthens mathematical problem-solving skills but also broadens the model’s overall reasoning abilities.

Following the pre-training phase, we perform mathematical instruction tuning on DeepSeekMath-Base using datasets that incorporate chain-of-thought \citep{cot}, program-of-thought \citep{pot,pal}, and tool-integrated reasoning \citep{tora} approaches. The model obtained, DeepSeekMath-Instruct 7B, outperforms all other models of similar (7B) scale and achieves results on par with open-source instruction-tuned models containing 70B parameters.

In addition, we propose Group Relative Policy Optimization (GRPO), a modified reinforcement learning (RL) algorithm based on Proximal Policy Optimization (PPO) \citep{schulman2017proximal}. Unlike standard approaches, GRPO eliminates the critic model and instead derives the baseline directly from group-level scores, thereby substantially cutting down on computational requirements for training. Utilizing only a limited selection of English instruction tuning data, GRPO delivers marked gains over the strong DeepSeekMath-Instruct model, achieving significant advances across both in-domain tasks (GSM8K: 82.9\% $\rightarrow$ 88.2\%, MATH: 46.8\% $\rightarrow$ 51.7\%) and out-of-domain benchmarks (e.g., CMATH: 84.6\% $\rightarrow$ 88.8\%) throughout RL training. Furthermore, we introduce a cohesive framework that positions various techniques—including Rejection Sampling Fine-Tuning (RFT) \citep{yuan2023scaling}, Direct Preference Optimization (DPO) \citep{dpo}, PPO, and GRPO—within a unified conceptual schema. Under this framework, these methodologies can be seen as either direct or streamlined forms of RL. We systematically explore the critical factors within this paradigm through comprehensive experimental studies—such as comparing online versus offline learning, evaluating outcome versus process-based supervision, and analyzing single-step versus iterative RL. Finally, we offer an explanation for the performance enhancements seen in instruction-tuned models under our RL regime, and we outline future avenues for developing even more effective RL strategies guided by this unified framework.

\subsection{Contributions}
We present advancements in scalable mathematics pre-training, complemented by an in-depth investigation and evaluation of reinforcement learning approaches.

\textbf{Math Pre-Training at Scale}

\begin{itemize}[topsep=0pt]
    \item This study demonstrates that the open Common Crawl dataset holds significant potential for mathematical applications. Through a carefully crafted pipeline for data selection, we build the DeepSeekMath Corpus—a robust dataset consisting of 120B tokens extracted from mathematically relevant web pages. This corpus is nearly sevenfold larger than the collection of math web pages utilized by Minerva \citep{minerva}, and exceeds the scale of OpenWebMath \citep{openwebmath} by a factor of nine.

\item DeepSeekMath-Base 7B, our foundational pre-trained model, demonstrates performance on par with Minerva 540B \citep{minerva}. This suggests that model size alone does not solely determine mathematical reasoning ability. In particular, a smaller model can exhibit robust performance if trained on sufficiently high-quality data.

\item We present results from our experiments involving math training. Exposing models to code training before math-focused training enhances their proficiency in solving mathematical problems, irrespective of tool assistance. This provides some insight into the ongoing debate: \textit{does code training enhance reasoning skills?} Our evidence indicates that it does, specifically with regard to mathematical reasoning.

\item While it is typical to use arXiv papers for training, particularly in numerous mathematics-focused studies, such a practice does not yield significant gains on any of the mathematical benchmarks considered in this work.
\end{itemize}

\textbf{Exploration and Analysis of Reinforcement Learning}

\begin{itemize}[topsep=0pt]
    \item We present Group Relative Policy Optimization (GRPO), a reinforcement learning method that prioritizes efficiency and effectiveness. Unlike traditional approaches that rely on a critic model, GRPO leverages group scores to estimate the baseline, thereby greatly decreasing the computational demands associated with Proximal Policy Optimization (PPO).
    \item Our findings indicate that applying GRPO substantially improves the performance of the instruction-tuned model DeepSeekMath-Instruct, utilizing only instruction-tuning datasets. In addition, we note performance gains on tasks outside the training domain as reinforcement learning progresses.
    \item We offer a cohesive framework for understanding various techniques, including RFT, DPO, PPO, and GRPO. Our work includes comprehensive experimental analyses—such as comparing online versus offline learning, examining the impact of outcome versus process supervision, and contrasting single-turn with iterative reinforcement learning—to thoroughly evaluate the critical aspects of this framework.
    \item Drawing from this unified framework, we investigate the mechanisms underlying the efficacy of reinforcement learning and highlight promising avenues for enhancing large language model reinforcement learning strategies.
\end{itemize}

\subsection{Summary of Evaluations and Metrics}

\begin{itemize}[topsep=0pt]
    \item \textbf{English and Chinese Mathematical Reasoning}: Our evaluation extensively measures model performance on both English and Chinese mathematical reasoning datasets, spanning from elementary to collegiate levels. The English datasets assessed comprise GSM8K \citep{gsm8k}, MATH \citep{MATH}, SAT \citep{llemma}, OCW Courses \citep{minerva}, and MMLU-STEM \citep{mmlu}. For Chinese, we utilize MGSM-zh \citep{mgsm}, CMATH \citep{wei2023cmath}, Gaokao-MathCloze \citep{agieval}, as well as Gaokao-MathQA \citep{agieval}. Our evaluation investigates both the capability of models to construct full-sentence solutions independently and their proficiency in resolving problems with the assistance of Python programming.

Across English evaluation tasks, DeepSeekMath-Base demonstrates performance on par with the proprietary Minerva 540B model \citep{minerva} and consistently outperforms all publicly available base models, such as Mistral 7B \citep{mistral} and Llemma-34B \citep{llemma}, regardless of their exposure to mathematical pre-training, often by a wide margin. Importantly, DeepSeekMath-Base achieves even stronger results on Chinese evaluation sets, which can be attributed to our inclusion of high-quality, non-English pre-training data—differing from previous studies \citep{minerva,llemma} that focused exclusively on English sources. After applying mathematical instruction tuning and reinforcement learning, DeepSeekMath-Instruct and DeepSeekMath-RL achieve impressive outcomes, becoming the first open-source models to surpass 50\% accuracy on the competition-level MATH benchmark.

\item \textbf{Formal Mathematics}: For formal mathematics evaluation, DeepSeekMath-Base is assessed on the informal-to-formal theorem proving benchmark introduced in \citep{dsp_proof}, utilizing miniF2F \citep{minif2f} with Isabelle \citep{isabelle} as the designated proof assistant. The results show that DeepSeekMath-Base excels in few-shot autoformalization tasks.

\item \textbf{Natural Language Understanding, Reasoning, and Code}: To thoroughly assess the model's proficiency in general comprehension, logical reasoning, and programming, we apply DeepSeekMath-Base to several standard benchmarks. This includes the Massive Multitask Language Understanding (MMLU) dataset \citep{mmlu} featuring 57 diverse multiple-choice questions, the BIG-Bench Hard (BBH) suite \citep{bbh} of 23 tasks emphasizing multi-step reasoning, as well as code-specific assessments with HumanEval \citep{codex} and MBPP \citep{mbpp}. Pre-training on mathematical corpora is found to enhance capabilities in both understanding and reasoning tasks.

\section{Math Pre-Training}

\subsection{Data Collection and Decontamination} This section details the methodology used to assemble the DeepSeekMath Corpus utilizing Common Crawl data. As illustrated in Figure \ref{fig:data_collect}, we introduce an iterative framework designed to methodically extract an extensive mathematical corpus from Common Crawl, beginning with an initial seed dataset comprised of a limited yet high-caliber selection of math-focused material. Notably, this methodology can be extended to additional fields, including but not limited to coding.

To begin, OpenWebMath \citep{openwebmath}, a repository containing curated mathematical web materials, is selected as the foundational seed corpus. Utilizing this dataset, we build a fastText model \citep{joulin2016fasttext} to retrieve additional web pages with content resembling that of OpenWebMath. For this process, we draw 500,000 random samples from the seed corpus to serve as positive instances and select an equal number of web pages from Common Crawl as negatives. The training procedure utilizes the official open-source fastText library\footnote{\url{https://fasttext.cc}}, with the parameters set to a 256-dimensional vector space, learning rate of 0.1, maximum word n-gram length of 3, minimum word occurrence threshold of 3, and 3 training epochs. To address the massive scale of Common Crawl, URL-based deduplication methods, along with near-duplicate filtering, are implemented, leaving a reduced set of approximately 40 billion HTML pages. The trained fastText classifier is then used to identify potential mathematical pages within this deduplicated dataset. To eliminate substandard content, pages are ranked based on fastText scores, and only those with the highest rankings are retained. The retained data volume is validated via pre-training studies that analyze datasets composed of the leading 40B, 80B, 120B, and 160B tokens. In this initial stage, only the top 40B tokens are selected for further use.

Following the initial round of data collection, a significant number of mathematical web pages are still missing from the dataset. This shortfall arises because the fastText model’s training set of positive samples is not sufficiently diverse. To remedy this, we augment the seed corpus by identifying further mathematical web resources, aiming to enhance the fastText model’s performance. Our procedure begins by partitioning the entire Common Crawl dataset into mutually exclusive domains, where a domain comprises web pages sharing the same base URL. For every domain, we determine the proportion of web pages captured during the first iteration. Domains with more than 10\% of their web pages collected are designated as math-oriented domains (for instance, \url{mathoverflow.net}). 
Within these domains, we proceed to manually label URLs that specifically correspond to mathematical content (such as \url{mathoverflow.net/questions}). Any web page associated with these labeled URLs that was missed in the initial round is subsequently incorporated into the seed corpus. This process secures a richer set of positive examples, allowing us to develop a more effective fastText model that can recover a larger share of mathematical web content in following rounds. Over four rounds of this iterative collection strategy, we accumulate a total of 35.5M mathematical web pages, amounting to 120B tokens. By the fourth iteration, we observe that approximately 98\% of the available data had already been acquired during the third round, leading us to terminate the collection effort at that point.

To prevent overlap with benchmark datasets, we adopt the approach outlined in \cite{deepseek-coder}, excluding web content that includes questions or solutions sourced from English mathematics benchmarks like GSM8K \citep{gsm8k} and MATH \citep{MATH}, as well as Chinese benchmarks such as CMATH \citep{wei2023cmath} and AGIEval \citep{agieval}. Our methodology for filtering is as follows: any portion of text in our math training data that contains a 10-gram sequence matching exactly with a corresponding substring from the benchmarks is discarded. In cases where benchmark materials consist of sequences shorter than 10 grams but at least 3 grams in length, we rely on direct string matches to identify and filter out web content that may result in data contamination.

\subsection{Validating the Quality of the DeepSeekMath~Corpus}

We run pre-training experiments to investigate how the DeepSeekMath~Corpus is compared with the recently released math-training corpora:

\begin{itemize}[topsep=0pt]
    \item \textbf{MathPile} \citep{mathpile}: This dataset compiles content from diverse sources including textbooks, Wikipedia, ProofWiki, CommonCrawl, StackExchange, and arXiv, comprising a total of 8.9B tokens. Notably, arXiv represents over 85\% of the dataset;
    \item \textbf{OpenWebMath} \citep{openwebmath}: Extracted from CommonCrawl by applying mathematical content filters, this collection amounts to 13.6B tokens;
    \item \textbf{Proof-Pile-2} \citep{llemma}: Constructed by merging OpenWebMath, AlgebraicStack (which contains 10.3B tokens derived from mathematical code), and arXiv documents (28.0B tokens in total). When evaluating on Proof-Pile-2, we adhere to the arXiv:Web:Code token ratio of 2:4:1 as established in \cite{llemma}.
\end{itemize}

\subsubsection{Training Setting}
\label{sec:quality-policy}
We fine-tune a general-purpose language model comprising 1.3B parameters—architecturally aligned with the DeepSeek LLMs \citep{deepseek-llm} and hence referred to as DeepSeek-LLM 1.3B—by subjecting it to mathematical training. For each mathematical dataset, an independent model is trained over a span of 150B tokens. All training processes utilize the HAI-LLM framework \citep{haillm}, known for its lightweight and resource-efficient nature. Adhering to the methodologies recommended by DeepSeek LLMs, we implement the AdamW optimizer \citep{adamW} configured with $\beta_1=0.9$, $\beta_2=0.95$, and $\mathrm{weight\_decay}=0.1$. The learning rate strategy involves an initial warmup for 2,000 steps until reaching its peak, followed by a reduction to 31.6\% after 80\% of training and a further decrease to 10.0\% of its maximum by the 90\% mark. The learning rate is capped at a maximum of 5.3e-4, and training employs a batch size of 4M tokens, utilizing a context window of 4K tokens.

\subsubsection{Evaluation Results}

\textbf{The DeepSeekMath~Corpus is of high quality, covers multilingual mathematical content, and is the largest in size.}

\begin{itemize}[topsep=0pt]
    \item \textbf{High-quality}:
    We assess downstream capabilities on eight mathematical tasks using few-shot chain-of-thought prompting as outlined in \cite{cot}. According to the results presented in Table \ref{tab:corpora_comparison}, models trained on the DeepSeekMath~Corpus consistently outperform counterparts trained on other datasets. Figure \ref{fig:corpora_comparisons} further illustrates that at the 50B-token point (equivalent to a complete epoch on Proof-Pile-2), the model utilizing the DeepSeekMath Corpus surpasses those trained on Proof-Pile-2, highlighting the superior average quality present in the DeepSeekMath Corpus.
    \item \textbf{Multilingual}:
    The DeepSeekMath~Corpus includes content in a variety of languages, with English and Chinese forming the bulk of the dataset. Table \ref{tab:corpora_comparison} demonstrates that training on DeepSeekMath~Corpus leads to significant advancements in mathematical reasoning in both English and Chinese. In comparison, other available mathematical corpora, which are heavily skewed towards English, either offer minimal gains or even degrade performance on mathematical reasoning tasks in Chinese.
    \item \textbf{Large-scale}:
    The scale of the DeepSeekMath~Corpus greatly exceeds that of previously available mathematical corpora. As shown in Figure \ref{fig:corpora_comparisons}, when DeepSeek-LLM 1.3B is trained on DeepSeekMath~Corpus, it exhibits a more rapid learning trajectory and more sustained progress. Conversely, baseline datasets are much smaller, have undergone repeated epochs of training, and models trained on them rapidly reach their optimal performance.
\end{itemize}

\subsection{Training and Evaluating DeepSeekMath-Base 7B}

In this section, we present DeepSeekMath-Base 7B, a foundational model designed for robust mathematical reasoning. The model initialization is based on DeepSeek-Coder-Base-v1.5 7B \citep{deepseek-coder} and it undergoes training over 500B tokens. The dataset composition is as follows: 56\% originates from the DeepSeekMath Corpus, 4\% is sourced from AlgebraicStack, 10\% comes from arXiv, 20\% consists of code from Github, and the final 10\% is natural language data, obtained from Common Crawl in both English and Chinese. For model training, we primarily utilize the setup detailed in Section \ref{sec:quality-policy}, with the modifications that the peak learning rate is increased to 4.2e-4 and the batch size is set at 10 million tokens.

We present an extensive evaluation of DeepSeekMath-Base 7B’s mathematical proficiency, examining its capacity to generate independent mathematical solutions, employ tools for problem-solving, and perform formal theorem verification. Additionally, our analysis extends to a broader characterization of the foundational model, encompassing its aptitude in natural language understanding, logical reasoning, and programming abilities.

\paragraph{Mathematical Problem Solving with Step-by-Step Reasoning}

We assess the capability of DeepSeekMath-Base in mathematical problem-solving under few-shot chain-of-thought prompting \citep{cot} across eight evaluation benchmarks, both in English and Chinese. The chosen benchmarks span a range of quantitative reasoning datasets—including GSM8K \citep{gsm8k}, MATH \citep{MATH}, and CMATH \citep{wei2023cmath}—as well as multiple-choice assessments such as MMLU-STEM \citep{mmlu} and Gaokao-MathQA \citep{agieval}. Together, these benchmarks represent a wide array of mathematical domains, from elementary concepts to college-level topics.

According to Table \ref{tab:base_math_cot}, DeepSeekMath-Base 7B demonstrates superior results across all eight evaluated benchmarks when compared with other open-source base models, such as the widely-adopted Mistral 7B \citep{mistral} and the recently introduced Llemma 34B \citep{llemma}, which has been specifically trained on Proof-Pile-2 \citep{llemma} for mathematical tasks. Particularly striking is its performance on the advanced MATH dataset, where DeepSeekMath-Base achieves an absolute margin of over 10\% beyond its open-source counterparts and even exceeds the accuracy of Minerva 540B \citep{minerva}. The latter, despite being 77 times larger and based on PaLM \citep{palm} with additional mathematical training, does not match the effectiveness of DeepSeekMath-Base.

\paragraph{Mathematical Problem Solving with Tool Use} We assess the capabilities of program-assisted mathematical reasoning on the GSM8K and MATH datasets through few-shot program-of-thought prompting \citep{pot,pal}. The models address each question by generating a Python script, making use of libraries like \textit{math} and \textit{sympy} to handle complex calculations. The final program output serves as the proposed solution. Table \ref{tab:base_math_pot_proof} indicates that DeepSeekMath-Base 7B achieves superior performance compared to the previous leading model, Llemma 34B.

\paragraph{Formal Mathematics}

The automation of formal proof has proven advantageous for guaranteeing the precision and dependability of mathematical proofs, as well as for improving productivity, drawing growing interest lately. In this study, we assess DeepSeekMath-Base 7B for its ability to convert informal mathematical content into formal proofs as described by \citep{dsp_proof}. Specifically, the model receives an informal statement, its formal version, and an informal proof, and must produce a corresponding formal proof. Our experiments are conducted on the miniF2F dataset \citep{minif2f}, a recognized benchmark featuring formalized Olympiad-level math problems. For each task, a formal proof in Isabelle is generated via few-shot prompting. Consistent with the methodology in \cite{dsp_proof}, we utilize the model to draft preliminary proof outlines, then apply the established automated tool Sledgehammer \citep{sledgehammer} to complete the finer points. Table \ref{tab:base_math_pot_proof} illustrates that DeepSeekMath-Base 7B achieves impressive results in automatic proof formalization.

\paragraph{Natural Language Understanding, Reasoning, and Code}

We assess the model’s abilities in natural language understanding using MMLU \citep{mmlu}, in reasoning using BBH \citep{bbh}, and in programming proficiency using HumanEval \citep{codex} and MBPP \citep{mbpp}. Table \ref{tab:understanding_reasoning_code} reveals that DeepSeekMath-Base 7B achieves considerable improvements on MMLU and BBH compared to its predecessor, DeepSeek-Coder-Base-v1.5 \citep{deepseek-coder}, highlighting the beneficial role of mathematical training in boosting both linguistic comprehension and reasoning. Furthermore, by incorporating code tokens throughout continual training, DeepSeekMath-Base 7B is able to preserve the coding performance of DeepSeek-Coder-Base-v1.5 on both coding datasets. In aggregate, DeepSeekMath-Base 7B demonstrates marked superiority over the general-purpose model Mistral 7B \citep{mistral} across the three evaluated reasoning and coding tasks.

\section{Supervised Fine-Tuning}

\subsection{SFT Data Curation}

We develop a dataset for instruction-tuning in mathematics that encompasses both English and Chinese questions drawn from a diverse array of mathematical domains and spanning multiple levels of difficulty. Each problem is accompanied by an answer formulated in either chain-of-thought (CoT) \citep{cot}, program-of-thought (PoT) \citep{pot,pal}, or tool-integrated reasoning \citep{tora} formats. Altogether, the dataset consists of 776K training instances.

\begin{itemize}[topsep=0pt]
    \item \textbf{English mathematical datasets}: We enrich GSM8K and MATH question sets with solutions that integrate computational tools, and further incorporate selected problems from MathInstruct \citep{MathInstruct} along with the training data from Lila-OOD \citep{lila}, all featuring solutions formulated through either CoT or PoT. Our English dataset encompasses a wide spectrum of mathematical areas, including algebra, probability, number theory, calculus, and geometry.
    \item \textbf{Chinese mathematical datasets}: We compile a collection of K-12 mathematics problems in Chinese, distributed across 76 distinct sub-domains such as linear equations, and provide both CoT-style and tool-based annotated solutions.
\end{itemize}

\subsection{Training and Evaluating DeepSeekMath-Instruct 7B}

This section presents DeepSeekMath-Instruct 7B, which is derived from DeepSeekMath-Base through mathematical instruction tuning. Training samples are combined in a random manner up to a total maximum context length of 4K tokens. The training procedure consists of 500 steps, utilizing a batch size of 256 and maintaining a fixed learning rate of 5e-5 throughout.

We evaluate models' mathematical performance both without and with tool use, on 4 quantitative reasoning benchmarks in English and Chinese. We benchmark our model against the leading models of the time:

\begin{itemize}[topsep=0pt]
    \item \textbf{Closed-source models} comprise: (1) the GPT series, with GPT-4 \citep{gpt4} and the GPT-4 Code Interpreter~\footnote{\url{https://openai.com/blog/chatgpt-plugins\#\#code-interpreter}} representing the most advanced, (2) Gemini Ultra and Gemini Pro \citep{gemini}, (3) Inflection-2 \citep{inflection-2}, (4) Grok-1~\footnote{\url{https://x.ai/model-card}}, in addition to newer systems from Chinese firms such as (5) Baichuan-3~\footnote{\url{https://www.baichuan-ai.com}}, and (6) the most recent GLM-4~\footnote{\url{https://open.bigmodel.cn/dev/api\#glm-4}}

    from the GLM family \citep{glm}.

These models are intended for broad use cases, with the majority having experienced multiple stages of alignment.
    \item \textbf{Open-source models} encompass: general-purpose models such as (1) DeepSeek-LLM-Chat 67B \citep{deepseek-llm}, (2) Qwen 72B \citep{qwen}, (3) SeaLLM-v2 7B \citep{seallm}, and (4) ChatGLM3 6B \citep{chatglm3}. In addition, several models provide specific advancements for mathematical tasks, including (5) InternLM2-Math 20B~\footnote{\url{https://github.com/InternLM/InternLM-Math}}, which extends InternLM2 via targeted mathematics training and subsequent instruction fine-tuning; (6) Math-Shepherd-Mistral 7B, which leverages PPO-based optimization \citep{schulman2017proximal} on the Mistral 7B framework \citep{mistral} coupled with a process-supervised reward strategy; (7) the WizardMath collection \citep{wizardmath}, aiming to bolster mathematical reasoning in Mistral 7B and Llama-2 70B \citep{llama2} by employing evolve-instruct—an approach to instruction tuning involving instructions generated by AI—as well as PPO training, using a data distribution predominantly from GSM8K and MATH; (8) MetaMath 70B \citep{metamath}, which refines Llama-2 70B on an enhanced dataset comprising GSM8K and MATH; (9) ToRA 34B \cite{tora}, based on CodeLlama 34B and further fine-tuned to support tool-augmented mathematical problem solving; and (10) MAmmoTH 70B \citep{MathInstruct}, which represents Llama-2 70B further aligned through MathInstruct-based instruction tuning.
\end{itemize}

Table \ref{tab:sft_rl_math} presents results for scenarios in which models are restricted from using external tools. In this setting, DeepSeekMath-Instruct 7B exhibits robust capability in multi-step reasoning tasks. Particularly on the MATH dataset, which targets competition-level problems, our model outperforms all existing open-source alternatives, as well as most closed-source models—including Inflection-2 and Gemini Pro—by a minimum margin of 9\% absolute. This superiority persists even against models with significantly larger parameter counts (e.g., Qwen 72B) and those trained with dedicated math-oriented reinforcement learning techniques (such as WizardMath-v1.1 7B). Although DeepSeekMath-Instruct achieves performance comparable to leading Chinese proprietary models like GLM-4 and Baichuan-3, it does not yet match the results obtained by GPT-4 or Gemini Ultra.

When assessed in a context that permits the combination of natural language reasoning with the utilization of programmatic tools for solving problems, DeepSeekMath-Instruct 7B achieves nearly 60\% accuracy on the MATH benchmark, outperforming every previously available open-source model. Across additional benchmarks, our model demonstrates performance on par with DeepSeek-LLM-Chat 67B, the former state-of-the-art despite its much greater size, being tenfold larger.

\subsection{Group Relative Policy Optimization} Reinforcement learning (RL) has demonstrated substantial success in enhancing the mathematical reasoning capabilities of LLMs following the Supervised Fine-Tuning (SFT) phase \citep{wang2023math,wizardmath}. Here, we present Group Relative Policy Optimization (GRPO), a novel RL algorithm designed to deliver both efficiency and effectiveness.

\subsubsection{From PPO to GRPO} Proximal Policy Optimization (PPO) \citep{schulman2017proximal} is an actor-critic RL algorithm that is widely used in the RL fine-tuning stage of LLMs \citep{ouyang2022training}. In particular, it optimizes LLMs by maximizing the following surrogate objective:

\begin{equation}
\footnotesize
    \mathcal{J}_{PPO}(\theta) = \mathbb{E}{[q \sim P(Q), o \sim \pi_{\theta_{old}}(O|q)]} \frac{1}{|o|} \sum_{t=1}^{|o|} \min \left[ \frac{\pi_\theta(o_{t} | q, o_{<t})}{\pi_{\theta_{old}}(o_{t} | q, o_{<t})} A_{t}, \text{clip} \left( \frac{\pi_\theta(o_{t} | q, o_{<t})}{\pi_{\theta_{old}}(o_{t} | q, o_{<t})}, 1 - \varepsilon, 1 + \varepsilon \right)  A_{t} \right] ,
\end{equation}

Here, $\pi_{\theta}$ refers to the current policy, while $\pi_{\theta_{old}}$ denotes the previous policy version. The variables $q$ and $o$ represent the questions sampled from the dataset and the outputs generated by the old policy $\pi_{\theta_{old}}$, respectively. The hyper-parameter $\varepsilon$ serves as a clipping threshold in PPO, playing a key role in making the training process more robust. The term $A_t$ signifies the advantage at time $t$, which is estimated via the Generalized Advantage Estimation (GAE) framework \citep{gae}, utilizing the rewards $\{r_{\ge t}\}$ and an associated value function $V_{\psi}$. Consequently, training in PPO requires jointly learning a value function in addition to the policy model. To prevent excessive optimization of the reward model, it is common practice to incorporate a KL divergence penalty (on a per-token basis) with respect to a reference model, added directly into the reward at each token position \citep{ouyang2022training}, that is,

\begin{equation}
    r_{t} = r_\varphi(q, o_{\le t}) - \beta \log\frac{\pi_{\theta}(o_{t}|q, o_{<t})}{\pi_{ref}(o_{t}|q, o_{<t})},
\label{eq:PPO-reward}
\end{equation}

Here, $r_\varphi$ denotes the reward model, while $\pi_{ref}$ refers to the reference model—typically, this is the original SFT model. The parameter $\beta$ represents the weight assigned to the KL penalty.

As the value function employed in PPO is typically another model of comparable size as the policy model, it brings a substantial memory and computational burden. Additionally, during RL training, the value function is treated as a baseline in the calculation of the advantage for variance reduction. While in the LLM context, usually only the last token is assigned a reward score by the reward model, which may complicate the training of a value function that is accurate at each token. To address this, as shown in Figure \ref{fig:grpo}, we propose Group Relative Policy Optimization (GRPO), which obviates the need for additional value function approximation as in PPO, and instead uses the average reward of multiple sampled outputs, produced in response to the same question, as the baseline. More specifically, for each question $q$, GRPO samples a group of outputs $\{o_1, o_2, \cdots, o_G\}$  from the old policy  $\pi_{\theta_{old}}$  and then optimizes the policy model by maximizing the following objective:

\begin{equation}
\footnotesize
\begin{split}
    \mathcal{J}_{GRPO}(\theta) &= \mathbb{E}{[q \sim P(Q), \{o_i\}_{i=1}^G \sim \pi_{\theta_{old}}(O|q)]}  \\
    & \frac{1}{G}\sum_{i=1}^G\frac{1}{|o_i|} \sum_{t=1}^{|o_i|} \Bigg( \min \left[ \frac{\pi_\theta(o_{i,t} | q, o_{i,<t})}{\pi_{\theta_{old}}(o_{i,t} | q, o_{i,<t})} \hat{A}_{i,t},\; \text{clip} \left( \frac{\pi_\theta(o_{i,t} | q, o_{i,<t})}{\pi_{\theta_{old}}(o_{i,t} | q, o_{i,<t})},\; 1 - \varepsilon, 1 + \varepsilon \right)  \hat{A}_{i,t} \right] \\
    &\hspace{7cm} - \beta \mathbb{D}_{KL}\left[\pi_{\theta} || \pi_{ref}\right] \Bigg) ,
\end{split}
\label{eq:GRPO-obj}
\end{equation}

where $\varepsilon$ and $\beta$ are hyper-parameters, and $\hat{A}_{i,t}$  is the advantage calculated based on relative rewards of the outputs inside each group only, which will be detailed in the following subsections. The group relative way that GRPO leverages to calculate the advantages, aligns well with the comparative nature of rewards models,  as reward models are typically trained on datasets of comparisons between outputs on the same question. Also note that, instead of adding KL penalty in the reward, GRPO regularizes by directly adding the KL divergence between the trained policy and the reference policy to the loss, avoiding complicating the calculation of  $\hat{A}_{i,t}$. And different from the KL penalty term used in (\ref{eq:PPO-reward}), we estimate the KL divergence with the following unbiased estimator \citep{kl_approx}:

\begin{equation}
\small
    \mathbb{D}_{KL}\left[\pi_{\theta} || \pi_{ref}\right] = \frac{\pi_{ref}(o_{i,t}|q,o_{i,<t})}{\pi_{\theta}(o_{i,t}|q,o_{i,<t})} - \log\left(\frac{\pi_{ref}(o_{i,t}|q,o_{i,<t})}{\pi_{\theta}(o_{i,t}|q,o_{i,<t})}\right) - 1,
\end{equation}

which is guaranteed to be positive.

\begin{algorithm}[t]
  \small
  \caption{Iterative Group Relative Policy Optimization}
  \textbf{Input} initial policy model $\pi_{\theta_{\text{init}}}$; reward models $r_\varphi$; task prompts $\mathcal{D}$; 
  hyperparameters $\varepsilon$, $\beta$, $\mu$
  \begin{algorithmic}[1]
    \State Initialize the policy model as $\pi_\theta \leftarrow \pi_{\theta_{\text{init}}}$
    \For{iteration = 1, \dots, I}
       \State Set reference model $\pi_{ref} \leftarrow \pi_{\theta}$
      \For{step = 1, \dots, M}
      \State Draw a mini-batch $\mathcal{D}_b$ from the dataset $\mathcal{D}$
      \State Assign $\pi_{\theta_{old}} \leftarrow \pi_{\theta}$ to preserve the previous policy
      \State For every question $q$ in $\mathcal{D}_b$, sample $G$ candidate responses $\{o_i\}_{i=1}^G \sim \pi_{\theta_{old}} (\cdot \mid q)$
      \State For each generated output $o_i$, determine reward values $\{r_i\}_{i=1}^{G}$ by evaluating via $r_{\varphi}$
      \State Calculate the group-relative advantage $\hat{A}_{i,t}$ for each token position $t$ in every $o_i$
      \For{GRPO iteration = 1, \dots, $\mu$}
        \State Refine $\pi_\theta$ by optimizing the GRPO criterion (Equation \ref{eq:GC-GRPO})
      \EndFor
    \EndFor 
    \State Update the reward model $r_\varphi$ by further training with data from a replay buffer. 
    \EndFor 
  \end{algorithmic}
  \textbf{Output} $\pi_\theta$
  \label{alg:iter-grpo}
\end{algorithm}

\subsubsection{Outcome Supervision RL with GRPO} Specifically, given a question $q$, a set of outputs $\{o_1, o_2, \cdots, o_G\}$ is generated by sampling from the previous iteration of the policy model $\pi_{\theta_{old}}$. Each output is evaluated with a reward model, resulting in $G$ associated reward values $\mathbf{r} = \{r_1, r_2, \cdots, r_G\}$. These rewards are subsequently standardized using z-score normalization, achieved by subtracting the mean of $\mathbf{r}$ and dividing by its standard deviation. With outcome supervision, each output $o_i$ receives a final, normalized reward, which is uniformly assigned as the advantage $\hat{A}_{i, t}$ for every token within $o_i$; concretely, $\hat{A}_{i, t} = \widetilde{r}_i = \frac{r_i - {\rm mean}(\mathbf{r})}{{\rm std}(\mathbf{r})}$. The policy is then trained to maximize the objective outlined in equation (\ref{eq:GRPO-obj}).

\subsubsection{Process Supervision RL with GRPO}

While outcome supervision offers a reward signal only upon completing the entire output, this may prove inadequate and inefficient when guiding the policy for intricate mathematical reasoning tasks. Inspired by \cite{wang2023math}, we investigate process supervision, which supplies feedback at each intermediate reasoning stage. More concretely, for a given question $q$ and $G$ sampled outputs denoted as $\{o_1, o_2, \cdots, o_G\}$, we leverage a process reward model to assign scores to every step in the outputs, leading to a set of rewards represented as $\mathbf{R} = \{ \{r_1^{index(1)},\cdots,r_1^{index(K_1)}\}, \cdots,  \{r_G^{index(1)},\cdots,r_G^{index(K_G)}\} \}$. Here, $index(j)$ designates the end token index corresponding to the $j$-th step, and $K_i$ refers to the total number of steps present in the $i$-th output. These reward values are standardized by centering and scaling with respect to the mean and standard deviation, specifically: $\widetilde{r}_i^{index(j)} = \frac{r_i^{index(j)} - {\rm mean(\mathbf{R})}}{{\rm std(\mathbf{R})}}$.

To proceed, the advantage assigned to each token is obtained by summing the normalized rewards of all subsequent steps, that is, $\hat{A}_{i, t} = \sum_{index(j) \ge t} \widetilde{r}_i^{index(j)}$. The policy is then optimized by maximizing the objective described in equation (\ref{eq:GRPO-obj}).

\subsubsection{Iterative RL with GRPO} During the reinforcement learning procedure, the previously trained reward model might become inadequate for effectively guiding the evolving policy model. To address this issue, we investigate iterative RL using GRPO. As illustrated in Algorithm \ref{alg:iter-grpo}, this method involves producing updated training data for the reward model by sampling from the latest policy model. The reward model is then updated iteratively through a replay strategy that maintains 10\% of past training data. Subsequently, the current policy model is adopted as the new reference model, and training of the policy model continues with supervision from the updated reward model.

\subsection{Training and Evaluating DeepSeekMath-RL}

Our reinforcement learning (RL) experiments are built upon the DeepSeekMath-Instruct 7B model. For RL training, we utilize approximately 144,000 questions in the chain-of-thought format, specifically those pertaining to GSM8K and MATH drawn from the SFT data. In order to examine how RL affects benchmarks that are not represented during RL training, we omit all other SFT questions. The reward model's training set is assembled in the same way as described by \citep{wang2023math}. Our initial reward model is trained with the DeepSeekMath-Base 7B architecture, using a learning rate of $2 \times 10^{-5}$. For the policy model under the GRPO algorithm, a learning rate of $1 \times 10^{-6}$ is selected, while the KL penalty coefficient is fixed at $0.04$. Each question is used to generate $64$ samples, with a maximum sequence length of $1024$, and a batch size of $1024$ during training. After every exploration phase, the policy model parameters are updated only once. In the evaluation phase, DeepSeekMath-RL 7B is assessed on the same benchmarks as DeepSeekMath-Instruct 7B. For DeepSeekMath-RL 7B, tasks involving GSM8K and MATH with chain-of-thought prompts are categorized as in-domain, whereas the other benchmarks serve as out-of-domain tasks.

Table \ref{tab:sft_rl_math} presents comparative results for both open- and closed-source models employing chain-of-thought and tool-integrated reasoning across English and Chinese evaluation benchmarks. The analysis reveals: 1) DeepSeekMath-RL 7B achieves accuracy rates of 88.2\% on GSM8K and 51.7\% on MATH using chain-of-thought reasoning, outperforming every open-source model within the 7B–70B parameter range and most closed-source alternatives. 2) Notably, DeepSeekMath-RL 7B’s training is restricted to chain-of-thought-formatted instruction data from GSM8K and MATH, initialized from DeepSeekMath-Instruct 7B. Even with this limited training dataset, DeepSeekMath-RL 7B consistently exceeds the performance of DeepSeekMath-Instruct 7B on all metrics, illustrating the strong benefits conferred by reinforcement learning.

\section{Discussion}
Here, we present an analysis of the results obtained from our pre-training and reinforcement learning (RL) experiments.

\subsection{Lessons Learnt in Pre-Training}

We begin by discussing our observations during the pre-training phase. Except where explicitly noted, our experiments follow the training configurations detailed in Section \ref{sec:quality-policy}. Importantly, all references to the DeepSeekMath Corpus within this section pertain to the 89B-token dataset collected during the second round of data gathering.

\subsubsection{Code Training Benefits Mathematical Reasoning}

An often-cited but largely untested conjecture posits that code training enhances reasoning abilities. In this work, we seek to address this claim, focusing specifically on mathematics: code training boosts models’ proficiency in mathematical reasoning tasks, regardless of whether external tools are used.

To study how code training affects mathematical reasoning, we experimented with the following two-stage training and one-stage training settings:

\textbf{Two-Stage Training}

\begin{itemize}[topsep=0pt]
    \item \textbf{Code Training for 400B Tokens $\rightarrow$ Math Training for 150B Tokens}: The DeepSeek-LLM 1.3B model undergoes initial pretraining on 400B code tokens, which is then succeeded by 150B tokens specifically focused on mathematical content;
    \item \textbf{General Training for 400B Tokens $\rightarrow$ Math Training for 150B Tokens}: For comparative purposes, we also conduct a control experiment wherein the model is first exposed to 400B general tokens, drawn from a comprehensive general-purpose dataset assembled by DeepSeek-AI, before being trained on 150B math tokens, aiming to discern the effect of using code tokens rather than general tokens for fostering mathematical reasoning abilities.
\end{itemize}

\textbf{One-Stage Training}

\begin{itemize}[topsep=0pt]
    \item \textbf{150B Math Tokens Training}: DeepSeek-LLM 1.3B is trained specifically on a dataset comprising 150B math-related tokens.
    \item \textbf{Joint Training with 400B Code Tokens and 150B Math Tokens}: Sequentially training with math tokens after code tokens has been observed to negatively impact coding abilities. To address this, we explore whether integrating code and math tokens within a single-phase training regime can concurrently enhance mathematical reasoning and mitigate issues such as catastrophic forgetting.
\end{itemize}

\paragraph{Results}
The downstream outcomes obtained from various training configurations are presented in Table \ref{tab:code-math} and Table \ref{tab:code-math-general-eval}.

Integrating code training supports the development of program-aided mathematical reasoning, regardless of whether a two-stage or one-stage training approach is employed. According to Table \ref{tab:code-math}, code training in the two-stage framework on its own brings notable gains in solving GSM8K and MATH tasks using Python, with additional advancement achieved after incorporating mathematical training in the subsequent phase. Notably, in the one-stage training paradigm, combining both code and mathematical tokens effectively reduces catastrophic forgetting—a challenge observed with the two-stage protocol—while also producing mutually beneficial effects for both coding tasks (Table \ref{tab:code-math-general-eval}) and program-based mathematical reasoning tasks (Table \ref{tab:code-math}).

Code training contributes to better mathematical reasoning skills even in the absence of tool assistance. Within the two-stage training framework, code training in the first stage alone yields noticeable improvements. Additionally, it facilitates a more effective second stage of math training, which collectively leads to optimal outcomes. In contrast, merging code tokens and math tokens within a single-stage training process appears to diminish mathematical reasoning capabilities without tool utilization. This may be attributable to DeepSeek-LLM 1.3B’s relatively small scale, which may not possess sufficient capacity to concurrently learn from both code and mathematical data during one-stage training.

\subsubsection{ArXiv Papers Seem Ineffective in Improving Mathematical Reasoning}

ArXiv papers are commonly included as a component of math pre-training data \citep{minerva,gpt-f,llemma,mathpile}. However, detailed analysis regarding their impact on mathematical reasoning has not been extensively conducted. Perhaps counter-intuitively, according to our experiments, arXiv papers seem ineffective in improving mathematical reasoning. We experiment with models of different sizes, including DeepSeek-LLM 1.3B and DeepSeek-Coder-Base-v1.5 7B \citep{deepseek-coder}, using arXiv corpora that underwent varied processing pipelines:

\begin{itemize}[topsep=0pt]
    \item \textbf{MathPile} \citep{mathpile}: a collection of 8.9 billion tokens constructed using specific data cleaning and filtering heuristics, where scientific papers from arXiv account for more than 85\% of the content;
    \item \textbf{ArXiv-RedPajama} \citep{redpajama}: comprises all arXiv LaTeX source files with elements such as preambles, comments, macros, and bibliographic entries stripped out, resulting in a dataset of 28.0 billion tokens.
\end{itemize}

For our experimental setup, DeepSeek-LLM 1.3B was trained for 150B tokens, and DeepSeek-Coder-Base-v1.5 7B was trained for 40B tokens, each utilizing a separate arXiv corpus. Results indicate that incorporating arXiv papers yields little to no benefit for enhancing mathematical reasoning abilities. When these models were trained exclusively on arXiv data, neither demonstrated significant gains; in fact, performance occasionally worsened across the suite of mathematical benchmarks assessed. This evaluation encompassed quantitative reasoning tasks such as GSM8K and MATH (see Table \ref{tab:arxiv-cot}), multiple-choice assessments like MMLU-STEM (Table \ref{tab:arxiv-cot}), as well as formal mathematics tests like miniF2F (Table \ref{tab:arxiv_atp}).

However, this conclusion has its limitations and should be taken with a grain of salt. We have not yet studied:

\begin{itemize}[topsep=0pt]
    \item The influence of arXiv tokens on particular mathematical tasks not explored in this study, for example, the task of transforming formal theorem statements or proofs into their informal equivalents;
    \item How arXiv tokens perform when integrated with additional data sources;
    \item If the advantages offered by arXiv papers persist or increase as model size is scaled up.
\end{itemize}

Thus, further exploration is required, which we leave for future studies.

\subsection{Insights of Reinforcement Learning}
\subsubsection{Moving Towards a Unified Paradigm} In this part, we introduce an integrated framework to systematically examine a range of training strategies, including SFT, RFT, DPO, PPO, and GRPO. Moreover, we perform empirical studies to investigate components of this overarching framework. Typically, the gradient of any given training approach with respect to the parameter $\theta$ can be formulated as follows:

\begin{equation}
    \nabla_{\theta}\mathcal{J}_{\textcolor{red}{\mathcal{A}}}(\theta) = \mathbb{E}[\underbrace{(q,o) \sim \textcolor{red}{\mathcal{D}}}_{Data \ Source}]\left( \frac{1}{|o|} \sum_{t=1}^{|o|}  \underbrace{GC_{{\mathcal{A}}}(q, o, t, \textcolor{red}{\pi_{{rf}}})}_{Gradient \ Coefficient}  \nabla_{\theta}\log \pi_{\theta}(o_t | q, o_{<t})\right).
\label{eq:objective}
\end{equation}

There exist three key components: 1) \textit{Data Source $\mathcal{D}$}, which determines the training data; 2) \textit{Reward Function $\pi_{{rf}}$}, which is the source of the training reward signal; 3) \textit{Algorithm $\mathcal{A}$}: which processes the training data and the reward signal to the gradient coefficient $GC$ that determines the magnitude of the penalty or reinforcement for the data. We analyze several representative methods based on such a unified paradigm:

\begin{itemize}
    \item  \textbf{Supervised Fine-tuning (SFT)}: In SFT, a pretrained model is adapted by training it further on a dataset curated by humans.
    \item \textbf{Rejection Sampling Fine-tuning (RFT)}: RFT builds upon SFT by additional training on a set of outputs generated by the SFT model in response to SFT prompts, retaining only those responses deemed correct for further fine-tuning.
    \item \textbf{Direct Preference Optimization (DPO)}: DPO enhances the SFT model by continuing training on newly generated responses from the SFT model, applying a pairwise DPO loss to optimize based on preference.
    \item \textbf{Online Rejection Sampling Fine-tuning (Online RFT)}: Unlike RFT, Online RFT initializes the policy with the SFT model, but the model is iteratively updated using outputs produced in real time by the current policy, fine-tuning on these freshly sampled augmentations.
    \item \textbf{PPO/GRPO}: PPO/GRPO start from the SFT-initialized policy model and proceed by training it through reinforcement learning, where the model is improved using outputs produced by its own evolving, real-time policy.
\end{itemize}

Table \ref{tab:data_policy} provides an overview of the constituent elements of these methods. A comprehensive explanation of the derivation steps can be found in Appendix \ref{app:analysis-rl}.

\paragraph{Observation about Data Source} We categorize data sources into two distinct types: online sampling and offline sampling. Online sampling refers to training data collected through the interactive rollouts of the current training policy, whereas offline sampling relies on data generated by the initial SFT model. Both RFT and DPO utilize offline sampling, in contrast to Online RFT and GRPO, which employ an online sampling approach.

Figure \ref{fig:rl-analysis} illustrates that Online RFT substantially surpasses RFT across both benchmarks. Initially, Online RFT performs similarly to RFT; however, as training advances, Online RFT demonstrates a pronounced performance advantage, highlighting the benefits of an online learning approach. This observation aligns with expectations, since at the start of training, the actor closely mirrors the SFT model and thus the sampled data shows only slight discrepancies. As training progresses and the actor diverges further, the resulting sampled data presents more substantial differences, and the capacity for real-time data sampling confers increasingly significant benefits.

\paragraph{Observation about Gradient Coefficient} The procedure by which the algorithm utilizes the reward signal involves converting it into a gradient coefficient, which in turn guides the modification of model parameters. In our experiments, the reward function is decomposed into two categories: `Rule' and `Model'. Here, `Rule' corresponds to an assessment of response quality premised on answer correctness, while `Model' involves employing a reward model—trained using judgments produced by the rule-based approach—to evaluate responses. As illustrated in Equations \ref{eq:GC-RFT} and \ref{eq:GC-GRPO}, a fundamental distinction exists between GRPO and Online RFT: specifically, GRPO exclusively modifies its gradient coefficient in accordance with the reward value assigned by the reward model. This mechanism enables GRPO to both encourage and discourage responses to varying extents, reflective of reward magnitude. By contrast, Online RFT does not employ such adaptive adjustment; it fails to penalize incorrect outputs and applies an identical reinforcement strength to all correct responses, regardless of their individual merit.

Figure \ref{fig:rl-analysis} illustrates that GRPO outperforms online RFT, underscoring the effectiveness of modifying positive and negative gradient coefficients. Moreover, GRPO+PS achieves better results than GRPO+OS, which demonstrates the advantages provided by step-sensitive, fine-grained gradient coefficients. We also investigate the impact of iterative RL by carrying out two iterations in our experimental setup. As presented in Figure \ref{fig:iter-rl}, the iterative RL approach yields marked enhancements in performance, with the most pronounced gains observed during the initial iteration.

\subsubsection{Why RL Works?} In this work, we apply reinforcement learning to a selected subset of instruction tuning datasets, observing notable improvements over the baseline instruction tuning models. To delve deeper into the efficacy of reinforcement learning, we assess both Pass@K and Maj@K metrics for the Instruct and RL variants across two evaluation sets. Figure \ref{fig:maj-pass} demonstrates that, while RL leads to higher Maj@K scores, it does not yield improvements in Pass@K. This suggests that the gains from RL are largely due to an increased likelihood of producing a correct answer among the top-$K$ outputs, indicating that RL primarily strengthens the robustness of the output distribution rather than advancing the core competencies of the model. In parallel, \citep{wang2023making} discuss a \textbf{misalignment problem} found in reasoning tasks for SFT models, revealing that reasoning abilities in such models can be augmented by adopting various preference alignment approaches \citep{yuan2023rrhf,song2023preference,wang2023making}.

\subsubsection{How to Achieve More Effective RL?}
\label{sec:effec_rl}
Our findings show that RL yields strong results on mathematical reasoning challenges. Furthermore, we introduce a cohesive framework that captures various prominent training strategies under a unified perspective. In this framework, these approaches are characterized as variations or abstractions of RL methodologies. According to Equation \ref{eq:objective}, the paradigm comprises three crucial elements: Data Source, Algorithm, and Reward Function. We outline several prospective avenues for further research focusing on these components.

\paragraph{Data Source} The foundation of all training approaches is the choice of data source. Within the framework of RL, we define the data source as the set of unlabeled questions, accompanied by outputs that are generated via sampling from the policy model. In this work, our experimental setup restricts the data source to those questions used during the instruction tuning phase, applying a simple nucleus sampling approach for output generation. We hypothesize that this design choice may partly explain why our RL pipeline leads to improvements exclusively in the Maj@K metric. Moving forward, we aim to evaluate our RL pipeline on question prompts that are out-of-distribution, potentially employing \textbf{advanced sampling (decoding) strategies}, such as methods inspired by tree-search algorithms \citep{yao2023tree}. Moreover, progress in \textbf{efficient inference techniques} \citep{xia-etal-2023-speculative,leviathan2023fast,kwon2023efficient,xia2024unlocking}, which critically impact the exploration capabilities of policy models, will also be of significant importance.

\paragraph{Algorithms} Algorithms utilize both data and reward signals to determine the gradient coefficient for updating model parameters. According to Equation \ref{eq:objective}, most contemporary methods fundamentally \textbf{TRUST} the reward function's signal, using it to increase or decrease the conditional probability assigned to specific tokens. Nevertheless, this approach cannot guarantee the reliability of the reward signal at all times, particularly in highly intricate tasks. For instance, the PRM800K datasets \citep{lightman2023let}, despite being meticulously labeled by skilled annotators, still exhibit roughly 20\% inaccurate annotations\footnote{\url{https://github.com/openai/prm800k/issues/12\#issuecomment-1728491852}}. In light of this, our investigation centers on reinforcement learning algorithms that are resilient to noisy reward signals. We anticipate that \textbf{WEAK-TO-STRONG} \citep{burns2023weak} alignment approaches may induce a transformative shift in learning algorithms.

\paragraph{Reward Function} The reward function provides the main source of learning signals during training. In RL, this function typically takes the form of a neural reward model. We identify three critical avenues for advancing reward models: 1) \textbf{Strategies to improve the generalization capacity of the reward model.} It is essential that the reward model generalizes well, accommodating both out-of-distribution queries and outputs produced by sophisticated decoding methods. Failing this, RL may only serve to reinforce the existing distribution of LLMs, rather than substantially enhancing their inherent abilities; 2) \textbf{Approaches to capture and quantify reward model uncertainty.} Incorporating uncertainty can potentially mediate the transition from weak reward models to algorithms supporting the weak-to-strong learning paradigm; 3) \textbf{Developing effective techniques for constructing high-quality process reward models} which supply nuanced, step-level feedback throughout reasoning processes \citep{lightman2023let,wang2023math}.

\section{Conclusion, Limitation, and Future Work} In this work, we introduce DeepSeekMath, a model that establishes state-of-the-art results among open-source systems on the competition-level MATH benchmark and closely rivals proprietary models in performance. DeepSeekMath~is based on the DeepSeek-Coder-v1.5 7B architecture, and it is further trained through continual pretraining on 500B tokens, which notably includes 120B mathematics-specific tokens extracted from Common Crawl. Our comprehensive ablation analysis demonstrates that web-based sources contain ample high-quality mathematical content, whereas the value of arXiv data may have been previously overestimated. We also propose Group Relative Policy Optimization (GRPO), an adaptation of Proximal Policy Optimization (PPO), which substantially enhances mathematical reasoning capability while reducing memory requirements. Experimental outcomes reveal that GRPO continues to be beneficial even when DeepSeekMath-Instruct 7B attains strong performance on established benchmarks. Furthermore, we present a unified framework for interpreting a collection of approaches, and outline several promising avenues to advance reinforcement learning efficiency.

While DeepSeekMath~demonstrates strong performance on tasks requiring quantitative reasoning, its proficiency in areas such as geometry and theorem-proving is not as robust compared to proprietary models. For example, our preliminary experiments reveal that the model struggles with questions involving triangles and ellipses, suggesting the existence of selection bias in the datasets used during pre-training and fine-tuning. Moreover, due to limitations in model size, DeepSeekMath~underperforms GPT-4 when it comes to leveraging few-shot learning. Unlike GPT-4, which benefits from additional examples in few-shot scenarios, DeepSeekMath~displays comparable outcomes in both zero-shot and few-shot settings. Moving forward, we intend to enhance our data selection mechanisms to assemble a higher-quality pre-training dataset. We also plan to investigate new strategies, as described in Section \ref{sec:effec_rl}, to achieve more efficient reinforcement learning for large language models.

\end{document}